\documentclass[11pt,a4paper]{report}
\usepackage[utf8]{inputenc}
\usepackage[T1]{fontenc}
\usepackage[french]{babel}
\usepackage{amsmath, amsfonts, amssymb}
\usepackage{tikz}
\usepackage{listings}
\usepackage{xcolor}
\usepackage{geometry}
\geometry{margin=2.5cm}

\lstset{
    language=Python,
    basicstyle=\ttfamily\small,
    keywordstyle=\color{blue},
    stringstyle=\color{red},
    commentstyle=\color{green!60!black},
    numbers=left,
    numberstyle=\tiny,
    stepnumber=1,
    numbersep=5pt,
    backgroundcolor=\color{gray!10},
    showspaces=false,
    showstringspaces=false,
    showtabs=false,
    frame=single,
    tabsize=4,
    captionpos=b,
    breaklines=true,
    breakatwhitespace=false,
    extendedchars=true,
    literate={é}{{\'e}}1 {è}{{\`e}}1 {à}{{\`a}}1 {ç}{{\c{c}}}1 {œ}{{\oe}}1 {ù}{{\`u}}1
}

\title{Formes lineaires, hyperplans, espace dual et orthogonalite en dimension finie}
\author{Charles EDOU NZE}
\date{\today}

\begin{document}

\maketitle
\tableofcontents

\chapter{Formes lineaires, hyperplans, espace dual et orthogonalite en dimension finie}

\section{Visualisation Geometrique (TikZ)}

\begin{figure}[h]
\centering
\begin{tikzpicture}[scale=1.5]
  % Espace 3D simulé
  \draw[thick,->] (0,0,0) -- (3,0,0) node[anchor=north east]{$x_1$};
  \draw[thick,->] (0,0,0) -- (0,3,0) node[anchor=north west]{$x_2$};
  \draw[thick,->] (0,0,0) -- (0,0,3) node[anchor=south]{$x_3$};

  % Hyperplan
  \filldraw[fill=blue!20, draw=blue!50!black, opacity=0.7] (-1,2,-1) -- (2,-1,-1) -- (3,0,2) -- (0,3,2) -- cycle;
  \node[blue!80!black] at (1,2,0) {$H = \ker \phi$};

  % Vecteur normal
  \draw[thick, red, ->] (1,1,0.5) -- (2,2,1.5) node[anchor=west]{$\phi$ (Vecteur Normal)};

  % Vecteur dans H
  \draw[thick, green!60!black, ->] (0.5,1,0.5) -- (1.5,0,0.5) node[anchor=west]{$v \in H$};
\end{tikzpicture}
\caption{Hyperplan vectoriel comme noyau d'une forme lineaire}
\end{figure}

\chapter{Jalon 11 : Formes linéaires, hyperplans, espace dual et orthogonalité en dimension finie}

\section{1. Présentation du concept clé}

La genèse de la notion de dualité puise ses racines dans l'observation fondamentale que tout objet géométrique ou algébrique gagne à être étudié non seulement de l'intérieur, par sa constitution atomique (ses coordonnées dans une base), mais aussi de l'extérieur, par la manière dont il interagit avec des instruments de mesure. Imaginez un espace tridimensionnel comme une salle emplie de particules. Un vecteur représente la position ou la vitesse d'une de ces particules. Une forme linéaire n'est pas un objet spatial de même nature ; c'est un détecteur, un regard posé sur l'espace. Si l'on conçoit la température comme variant linéairement dans la pièce, le thermomètre évalue chaque position et lui attribue un scalaire unique. L'ensemble de ces "regards" possibles, de ces instruments d'évaluation, forme un nouvel univers : l'espace dual.

L'hyperplan, dans cette optique, s'impose comme la frontière naturelle. En deux dimensions, c'est la ligne parfaite tranchant le plan ; en trois dimensions, le feuillet infiniment mince scindant l'espace. Structurellement, l'hyperplan est le lieu géométrique de silence absolu pour une forme linéaire donnée, c'est-à-dire l'ensemble exact des vecteurs que l'instrument de mesure évalue à zéro, le "niveau de la mer" de notre topographie abstraite. C'est cette interaction entre l'objet mesuré (le primal) et l'instrument de mesure (le dual) qui donne naissance à la notion profonde d'orthogonalité, généralisant la perpendicularité euclidienne au-delà des structures dotées d'un produit scalaire usuel.

\section{2. Formalisation}

\subsection{A. Définitions Formelles et Typage Rigoureux}

Soit $E$ un espace vectoriel sur un corps commutatif $\mathbb{K}$ (typiquement $\mathbb{R}$ ou $\mathbb{C}$). On suppose dans la suite que la dimension de $E$ est finie, notée $n \in \mathbb{N}^*$.

\textbf{Forme linéaire :}
Une forme linéaire sur $E$ est une application $\phi : E \to \mathbb{K}$ qui est strictement linéaire sur le corps $\mathbb{K}$. Formellement, on requiert :
$$\forall x, y \in E, \forall \lambda, \mu \in \mathbb{K}, \quad \phi(\lambda x + \mu y) = \lambda \phi(x) + \mu \phi(y)$$

\textbf{Espace Dual ($E^*$) :}
L'espace dual de $E$, noté $E^*$, est l'espace vectoriel constitué de l'ensemble de toutes les formes linéaires sur $E$. Il s'identifie à l'espace des applications linéaires $\mathcal{L}(E, \mathbb{K})$. L'addition et la multiplication par un scalaire dans $E^*$ sont héritées point par point :
$$(\phi_1 + \phi_2)(x) = \phi_1(x) + \phi_2(x) \quad \text{et} \quad (\lambda \phi)(x) = \lambda \phi(x)$$

\textbf{Hyperplan :}
Un sous-espace vectoriel $H$ de $E$ est un hyperplan si et seulement s'il existe une forme linéaire non nulle $\phi \in E^*$ (soit $\phi \neq 0_{E^*}$) telle que $H = \ker \phi$. Le noyau de $\phi$, défini par $\{x \in E \mid \phi(x) = 0_\mathbb{K}\}$, caractérise entièrement cet espace frontière.

\textbf{Base Duale :}
Soit $\mathcal{B} = (e_1, \dots, e_n)$ une base de $E$. On définit la famille d'applications $\mathcal{B}^* = (e^*_1, \dots, e^*_n)$ de $E^*$ par leurs évaluations sur les vecteurs de base :
$$e^*_i(e_j) = \delta_{i,j}$$
où $\delta_{i,j}$ est le symbole de Kronecker (valant $1$ si $i=j$ et $0$ sinon). Par prolongement linéaire, $e^*_i(x)$ renvoie la $i$-ème coordonnée de $x$ dans la base $\mathcal{B}$. La famille $\mathcal{B}^*$ constitue une base de l'espace dual $E^*$, appelée base duale.

\textbf{Orthogonalité (Dualité) :}
Pour toute partie $A \subseteq E$, on définit l'orthogonal de $A$ (dans le dual) comme l'ensemble des annihilateurs de $A$ :
$$A^\perp = \{ \phi \in E^* \mid \forall x \in A, \phi(x) = 0 \}$$
Symétriquement, pour $B \subseteq E^*$, on définit $B^\circ = \{ x \in E \mid \forall \phi \in B, \phi(x) = 0 \}$.

\subsection{B. Théorèmes Fondamentaux}

\textbf{Théorème de la dimension du Dual :}
Si $\dim(E) = n$, alors l'espace dual $E^*$ est de dimension finie et $\dim(E^*) = n$.
*Corollaire structural :* Si $H$ est un hyperplan de $E$, alors $\dim(H) = n - 1$.

\textbf{Théorème du Bidual et Isomorphisme Canonique :}
Soit $E^{\textbf{} = (E^*)^*$ l'espace bidual. L'application d'évaluation $\Psi : E \to E^{}}$ définie par $\Psi(x)(\phi) = \phi(x)$ pour tout $x \in E$ et $\phi \in E^*$ est un isomorphisme canonique (indépendant du choix d'une base) dès lors que $E$ est de dimension finie. On a ainsi $E \cong E^{**}$.

\subsection{C. Exemples et Cas Pathologiques}

\textbf{Exemple de validation (Formes coordonnées) :}
Dans $\mathbb{R}^3$, pour un vecteur $x = (x_1, x_2, x_3)$, l'application $\phi(x) = 2x_1 - 3x_2 + x_3$ est une forme linéaire. Son noyau définit l'hyperplan vectoriel d'équation $2x_1 - 3x_2 + x_3 = 0$, qui est un plan passant par l'origine.

\textbf{Cas pathologique (Dimension infinie) :}
Si $E = \mathbb{R}[X]$ (l'espace des polynômes, de dimension infinie), la famille $(X^k)_{k \in \mathbb{N}}$ est une base. On peut définir les formes coordonnées $e_k^*$ par $e_k^*(X^j) = \delta_{k,j}$. Toutefois, la famille $(e_k^*)_{k \in \mathbb{N}}$ n'engendre pas $E^*$ ! En effet, la forme linéaire $\phi(P) = P(1) = \sum a_k$ ne peut s'écrire comme une combinaison linéaire finie des $e_k^*$. En dimension infinie, $E$ et $E^*$ ne sont pas isomorphes ($E^*$ a une dimension strictement "plus grande" que $E$).

\section{3. Démonstrations}

\subsection{Démonstration de la Dimension d'un Hyperplan}

Soit $H$ un hyperplan de $E$, un $\mathbb{K}$-espace vectoriel de dimension $n$. Montrons que $\dim(H) = n - 1$.

1. Par la définition formelle d'un hyperplan, il existe une forme linéaire $\phi \in E^*$ telle que $\phi$ n'est pas l'application nulle ($\phi \neq 0_{E^*}$) et $H = \ker \phi$.
2. Considérons l'application linéaire $\phi : E \to \mathbb{K}$. Le théorème du rang nous assure l'égalité structurelle :
   $$\dim(E) = \dim(\ker \phi) + \text{rg}(\phi)$$
   En substituant $H$ et $n$, nous obtenons $n = \dim(H) + \dim(\text{Im }\phi)$.
3. Le sous-espace $\text{Im }\phi$ est un sous-espace vectoriel du corps $\mathbb{K}$ (considéré comme $\mathbb{K}$-espace vectoriel de dimension $1$). Les seuls sous-espaces de $\mathbb{K}$ sont $\{0_\mathbb{K}\}$ et $\mathbb{K}$.
4. Puisque la forme linéaire $\phi$ n'est pas identiquement nulle, il existe nécessairement au moins un vecteur $v \in E$ tel que $\phi(v) \neq 0_\mathbb{K}$. L'image $\text{Im }\phi$ n'est donc pas réduite à $\{0_\mathbb{K}\}$.
5. Conséquemment, $\text{Im }\phi = \mathbb{K}$, ce qui implique que $\dim(\text{Im }\phi) = \dim(\mathbb{K}) = 1$. $\phi$ est donc surjective.
6. L'équation issue du théorème du rang devient :
   $$n = \dim(H) + 1$$
7. En isolant la dimension de l'hyperplan, on déduit irréfutablement :
   $$\dim(H) = n - 1$$

\subsection{Démonstration de l'Isomorphisme Canonique vers le Bidual}

Soit $\Psi : E \to E^{**}$ l'application définie pour tout $x \in E$ par $\Psi(x) = \text{ev}_x$, où $\text{ev}_x(\phi) = \phi(x)$ pour toute forme $\phi \in E^*$. Montrons que $\Psi$ est un isomorphisme lorsque $\dim(E) = n$.

1. \textbf{Vérification de la nature de l'application :}
   Pour tout $x \in E$, $\Psi(x)$ est bien une forme linéaire sur $E^*$. En effet, pour $\phi_1, \phi_2 \in E^*$ et $\alpha, \beta \in \mathbb{K}$ :
   $$\Psi(x)(\alpha \phi_1 + \beta \phi_2) = (\alpha \phi_1 + \beta \phi_2)(x) = \alpha \phi_1(x) + \beta \phi_2(x) = \alpha \Psi(x)(\phi_1) + \beta \Psi(x)(\phi_2)$$
   Ainsi, $\Psi(x) \in (E^*)^* = E^{**}$.

2. \textbf{Linéarité de $\Psi$ :}
   Montrons que $\Psi$ est elle-même linéaire. Soient $x, y \in E$ et $\lambda, \mu \in \mathbb{K}$. Il faut démontrer que $\Psi(\lambda x + \mu y) = \lambda \Psi(x) + \mu \Psi(y)$ dans $E^{**}$.
   Évaluons ces deux membres sur une forme quelconque $\phi \in E^*$ :
   $$\Psi(\lambda x + \mu y)(\phi) = \phi(\lambda x + \mu y)$$
   Par la stricte linéarité de $\phi$ sur $E$, nous déployons :
   $$\phi(\lambda x + \mu y) = \lambda \phi(x) + \mu \phi(y) = \lambda \Psi(x)(\phi) + \mu \Psi(y)(\phi) = (\lambda \Psi(x) + \mu \Psi(y))(\phi)$$
   Cette égalité étant vraie pour toute forme $\phi \in E^*$, on conclut que $\Psi(\lambda x + \mu y) = \lambda \Psi(x) + \mu \Psi(y)$. L'application $\Psi$ est donc linéaire.

3. \textbf{Injectivité de $\Psi$ :}
   Caractérisons le noyau de $\Psi$. Soit $x \in \ker \Psi$. Cela signifie que $\Psi(x) = 0_{E^{**}}$, soit pour toute forme linéaire $\phi \in E^*$, $\phi(x) = 0_\mathbb{K}$.
   Supposons par l'absurde que $x \neq 0_E$. Comme $E$ est de dimension finie $n$, et $x$ est un vecteur non nul, le théorème de la base incomplète nous autorise à construire une base de $E$ de la forme $\mathcal{B} = (e_1, e_2, \dots, e_n)$ où le premier vecteur est posé comme $e_1 = x$.
   Considérons la base duale $\mathcal{B}^* = (e_1^*, e_2^*, \dots, e_n^*)$. Par définition, la forme linéaire $e_1^*$ vérifie $e_1^*(e_1) = 1_\mathbb{K}$.
   Mais puisque $e_1 = x$, nous aurions $e_1^*(x) = 1_\mathbb{K}$.
   Or, l'hypothèse $x \in \ker \Psi$ impose que $\phi(x) = 0_\mathbb{K}$ pour absolument toute forme $\phi$, y compris $e_1^*$. Nous obtenons la contradiction $1_\mathbb{K} = 0_\mathbb{K}$.
   L'hypothèse $x \neq 0_E$ est donc erronée. Le noyau est réduit au vecteur nul : $\ker \Psi = \{0_E\}$. L'application $\Psi$ est strictement injective.

4. \textbf{Conclusion par argument de dimension :}
   Nous savons que $\dim(E^*) = \dim(E) = n$. Par application successive du même théorème, la dimension de l'espace bidual est $\dim(E^{**}) = \dim(E^*) = n$.
   Puisque $\Psi : E \to E^{**}$ est une application linéaire injective entre deux espaces vectoriels de même dimension finie $n$, le théorème de l'isomorphisme en dimension finie stipule que $\Psi$ est nécessairement surjective, et donc bijective. C'est un isomorphisme canonique.

\section{4. Exercices d'Application}

\subsection{Exercice 1 : Explicitation d'une Base Duale}
\textbf{Énoncé :}
Soit $E = \mathbb{R}^3$. On considère la famille de vecteurs $\mathcal{B} = (v_1, v_2, v_3)$ avec $v_1 = (1, 1, 0)$, $v_2 = (0, 1, 1)$ et $v_3 = (1, 0, 1)$. Démontrer que $\mathcal{B}$ est une base de $\mathbb{R}^3$ et construire explicitement la base duale associée $\mathcal{B}^* = (v_1^*, v_2^*, v_3^*)$ exprimée en fonction des coordonnées canoniques $(x, y, z)$.

\textbf{Correction exhaustive :}
1. \textbf{Démonstration de la nature de base de $\mathcal{B}$ :}
   Il suffit de montrer que la famille est libre. Posons $\alpha_1 v_1 + \alpha_2 v_2 + \alpha_3 v_3 = 0$.
   $\alpha_1(1, 1, 0) + \alpha_2(0, 1, 1) + \alpha_3(1, 0, 1) = (0, 0, 0)$
   Cela induit le système :
   (i) $\alpha_1 + \alpha_3 = 0$
   (ii) $\alpha_1 + \alpha_2 = 0$
   (iii) $\alpha_2 + \alpha_3 = 0$
   De (i), $\alpha_3 = -\alpha_1$. De (ii), $\alpha_2 = -\alpha_1$. En remplaçant dans (iii) : $-\alpha_1 - \alpha_1 = 0 \implies -2\alpha_1 = 0 \implies \alpha_1 = 0$.
   Par conséquent, $\alpha_1 = \alpha_2 = \alpha_3 = 0$. La famille est libre, et possédant 3 vecteurs en dimension 3, elle constitue une base de $E$.

2. \textbf{Construction de $v_1^*$ :}
   Cherchons $v_1^*$ sous la forme $v_1^*(x, y, z) = ax + by + cz$. Les relations de dualité exigent :
   - $v_1^*(v_1) = 1 \implies a + b = 1$
   - $v_1^*(v_2) = 0 \implies b + c = 0 \implies c = -b$
   - $v_1^*(v_3) = 0 \implies a + c = 0 \implies a = -c = b$
   En substituant $a = b$ dans la première équation : $b + b = 1 \implies b = \frac{1}{2}$. On tire $a = \frac{1}{2}$ et $c = -\frac{1}{2}$.
   Ainsi, $v_1^*(x, y, z) = \frac{1}{2}x + \frac{1}{2}y - \frac{1}{2}z$.

3. \textbf{Construction de $v_2^*$ :}
   Cherchons $v_2^*(x, y, z) = a'x + b'y + c'z$.
   - $v_2^*(v_1) = 0 \implies a' + b' = 0 \implies a' = -b'$
   - $v_2^*(v_2) = 1 \implies b' + c' = 1$
   - $v_2^*(v_3) = 0 \implies a' + c' = 0 \implies c' = -a' = b'$
   Dans la deuxième : $b' + b' = 1 \implies b' = \frac{1}{2}$. D'où $c' = \frac{1}{2}$ et $a' = -\frac{1}{2}$.
   Ainsi, $v_2^*(x, y, z) = -\frac{1}{2}x + \frac{1}{2}y + \frac{1}{2}z$.

4. \textbf{Construction de $v_3^*$ :}
   Cherchons $v_3^*(x, y, z) = a''x + b''y + c''z$.
   - $v_3^*(v_1) = 0 \implies a'' + b'' = 0 \implies b'' = -a''$
   - $v_3^*(v_2) = 0 \implies b'' + c'' = 0 \implies c'' = -b'' = a''$
   - $v_3^*(v_3) = 1 \implies a'' + c'' = 1 \implies a'' + a'' = 1 \implies a'' = \frac{1}{2}$.
   D'où $c'' = \frac{1}{2}$ et $b'' = -\frac{1}{2}$.
   Ainsi, $v_3^*(x, y, z) = \frac{1}{2}x - \frac{1}{2}y + \frac{1}{2}z$.

\subsection{Exercice 2 : Orthogonal d'un Sous-espace}
\textbf{Énoncé :}
Soit $E$ un espace vectoriel de dimension $n$. Montrer que si $F$ et $G$ sont deux sous-espaces de $E$ tels que $F \subseteq G$, alors au sens de la dualité, $G^\perp \subseteq F^\perp$. Montrer ensuite que $\dim(F) + \dim(F^\perp) = n$.

\textbf{Correction exhaustive :}
1. \textbf{Inclusion des orthogonaux :}
   Soit $\phi \in G^\perp$. Par la définition de l'orthogonal dans le dual, cela signifie que pour tout vecteur $x \in G$, nous avons $\phi(x) = 0$.
   Puisque l'hypothèse garantit que $F \subseteq G$, tout vecteur $y \in F$ est aussi un élément de $G$.
   Par conséquent, l'évaluation de $\phi$ sur tout vecteur de $F$ donne $\phi(y) = 0$.
   Ceci certifie que $\phi$ annule entièrement $F$. Ainsi, $\phi \in F^\perp$. L'inclusion $G^\perp \subseteq F^\perp$ est rigoureusement prouvée.

2. \textbf{Équation de dimension (Théorème d'isomorphisme canonique) :}
   Considérons l'application de restriction $\rho : E^* \to F^*$ définie par $\rho(\phi) = \phi_{|F}$, qui à une forme linéaire sur $E$ associe sa restriction au sous-espace $F$.
   L'application $\rho$ est clairement linéaire.
   Identifions son noyau. Une forme $\phi$ appartient à $\ker \rho$ si et seulement si sa restriction à $F$ est l'application nulle, c'est-à-dire que pour tout $x \in F, \phi(x) = 0$. C'est l'exacte définition de $F^\perp$. Donc $\ker \rho = F^\perp$.
   Démontrons la surjectivité de $\rho$. Soit $\psi \in F^*$ une forme linéaire définie uniquement sur $F$. Soit un supplémentaire $S$ de $F$ dans $E$ (tel que $E = F \oplus S$). Tout vecteur $x \in E$ se décompose de manière unique en $x = y + z$ avec $y \in F, z \in S$. On définit $\tilde{\psi} \in E^*$ par $\tilde{\psi}(x) = \psi(y) + 0$. Cette forme $\tilde{\psi}$ prolonge $\psi$ à l'espace total $E$, ce qui prouve que toute forme de $F^*$ admet un antécédent. L'application $\rho$ est donc surjective, et $\text{Im }\rho = F^*$.
   Par l'application rigoureuse du théorème du rang à l'application $\rho$ :
   $$\dim(E^*) = \dim(\ker \rho) + \dim(\text{Im }\rho)$$
   Nous savons que $\dim(E^*) = n$ et $\dim(F^*) = \dim(F)$ (puisque la dimension du dual est égale à celle du primal).
   En substituant : $n = \dim(F^\perp) + \dim(F)$.

\section{5. Application en Intelligence Artificielle}

Dans le cadre du Machine Learning et particulièrement de l'apprentissage statistique géométrique, la notion d'hyperplan est le socle de l'algorithme des machines à vecteurs de support (Support Vector Machines, SVM). L'objectif est de trouver un classifieur linéaire capable de discriminer deux catégories de données (par exemple des vecteurs $x_i \in \mathbb{R}^n$ étiquetés $+1$ et $-1$). Ce séparateur optimal est géométriquement décrit comme un hyperplan, dont l'équation caractéristique est modélisée par une forme linéaire $\phi(x) + b = 0$.
La puissance de la modélisation réside dans l'exploitation du théorème de bidualité. La minimisation de la norme du classifieur (le problème primal) est transformée, via le Lagrangien, en un problème de maximisation sur l'espace dual. Dans ce cadre dual, les variables d'optimisation (les multiplicateurs de Lagrange) agissent exactement comme des composantes de formes linéaires évaluant le poids critique (le "support") de chaque donnée d'apprentissage, réduisant ainsi drastiquement la complexité calculatoire pour les espaces de très haute dimension.

\section{6. Liens Sémantiques}
- \textbf{Concepts Précédents requis :} [[Jalon 7 (Espaces vectoriels abstraits)]], [[Jalon 8 (Applications linéaires)]]
- \textbf{Concepts Futurs dépendants :} [[Jalon 12 (Livrable IA)]], [[Jalon 25 (Formes bilinéaires)]], [[Jalon 123 (Problèmes d'optimisation sous contraintes)]]

\chapter{Exercices d'application}
\section{Exercice 1: Espace dual en dimension 2 (Difficulté $\star$)}

\subsection{Énoncé}
Soit $E = \mathbb{R}^2$. On considère la base canonique $(e_1, e_2)$. Déterminer la base duale $(e_1^*, e_2^*)$ et calculer l'évaluation de la forme linéaire $e_1^*$ sur le vecteur $v = 3e_1 - 2e_2$.

\subsection{Correction détaillée}

1. \textbf{Définition de la base duale :}
   L'espace $E = \mathbb{R}^2$ étant de dimension 2, son dual $E^*$ est également de dimension 2. La base duale associée à $(e_1, e_2)$ est la famille $(e_1^*, e_2^*)$ d'applications linéaires de $\mathbb{R}^2$ vers $\mathbb{R}$ définie explicitement par l'action sur les vecteurs de base :
   - $e_1^*(e_1) = 1$ et $e_1^*(e_2) = 0$
   - $e_2^*(e_1) = 0$ et $e_2^*(e_2) = 1$

2. \textbf{Explicitation des formes linéaires sur un vecteur quelconque :}
   Soit un vecteur $u = xe_1 + ye_2 \in \mathbb{R}^2$. Par la stricte linéarité de l'application $e_1^*$, nous pouvons développer l'évaluation :
   $$e_1^*(u) = e_1^*(xe_1 + ye_2) = x e_1^*(e_1) + y e_1^*(e_2)$$
   En substituant les valeurs des évaluations sur la base :
   $$e_1^*(u) = x \cdot 1 + y \cdot 0 = x$$
   De manière rigoureusement symétrique pour $e_2^*$ :
   $$e_2^*(u) = e_2^*(xe_1 + ye_2) = x e_2^*(e_1) + y e_2^*(e_2) = x \cdot 0 + y \cdot 1 = y$$
   Les éléments de la base duale sont donc les applications d'extraction des coordonnées canoniques.

3. \textbf{Calcul de l'évaluation spécifique :}
   Considérons le vecteur $v = 3e_1 - 2e_2$. Nous appliquons la linéarité de $e_1^*$ :
   $$e_1^*(3e_1 - 2e_2) = 3 e_1^*(e_1) - 2 e_1^*(e_2)$$
   En remplaçant par les définitions :
   $$e_1^*(3e_1 - 2e_2) = 3 \times 1 - 2 \times 0 = 3$$

\textbf{Conclusion :}
La valeur de l'évaluation est $3$. L'action de la forme coordonnée $e_1^*$ a parfaitement isolé la composante du vecteur $v$ sur le premier axe de base.

\section{Exercice 2: Caractérisation d'un hyperplan vectoriel (Difficulté $\star \star$)}

\subsection{Énoncé}
Dans $E = \mathbb{R}^3$, on considère le sous-espace $H = \{ (x, y, z) \in \mathbb{R}^3 \mid 2x - y + 3z = 0 \}$. Démontrer rigoureusement que $H$ est un hyperplan et expliciter une forme linéaire l'engendrant en tant que noyau.

\subsection{Correction détaillée}

1. \textbf{Définition de l'application candidate :}
   Considérons l'application $\phi : \mathbb{R}^3 \to \mathbb{R}$ définie par $\phi(x, y, z) = 2x - y + 3z$.

2. \textbf{Démonstration de la linéarité de $\phi$ :}
   Soient $u = (x_1, y_1, z_1)$ et $v = (x_2, y_2, z_2)$ deux vecteurs de $\mathbb{R}^3$. Soient $\alpha, \beta \in \mathbb{R}$.
   $$\phi(\alpha u + \beta v) = \phi(\alpha x_1 + \beta x_2, \alpha y_1 + \beta y_2, \alpha z_1 + \beta z_2)$$
   $$= 2(\alpha x_1 + \beta x_2) - (\alpha y_1 + \beta y_2) + 3(\alpha z_1 + \beta z_2)$$
   En réarrangeant les termes scalaires :
   $$= \alpha(2x_1 - y_1 + 3z_1) + \beta(2x_2 - y_2 + 3z_2)$$
   $$= \alpha \phi(u) + \beta \phi(v)$$
   L'application $\phi$ est donc une forme linéaire (un élément de $E^*$).

3. \textbf{Vérification de la non-nullité :}
   Pour s'assurer que $\phi$ engendre un hyperplan (et non l'espace total), il faut que $\phi$ ne soit pas l'application nulle.
   Évaluons $\phi$ sur le vecteur de la base canonique $e_1 = (1, 0, 0)$ :
   $$\phi(1, 0, 0) = 2(1) - 0 + 3(0) = 2 \neq 0$$
   Donc $\phi \neq 0_{E^*}$.

4. \textbf{Lien avec le noyau :}
   Par définition, l'ensemble $H$ est exactement constitué des vecteurs $u \in \mathbb{R}^3$ tels que $\phi(u) = 0$.
   Ainsi, $H = \ker \phi$.

\textbf{Conclusion :}
L'ensemble $H$ étant le noyau d'une forme linéaire non identiquement nulle sur un espace de dimension $3$, $H$ est formellement un hyperplan de $\mathbb{R}^3$, et par le théorème des dimensions, $\dim(H) = 3 - 1 = 2$.

\section{Exercice 3: Base duale d'une famille de polynômes (Difficulté $\star \star \star$)}

\subsection{Énoncé}
Dans l'espace vectoriel $E = \mathbb{R}_2[X]$ des polynômes de degré inférieur ou égal à 2, on considère la famille $\mathcal{B} = (P_0, P_1, P_2)$ où $P_0(X) = 1$, $P_1(X) = X-1$, et $P_2(X) = (X-1)^2$. Montrer que $\mathcal{B}$ est une base et déterminer la base duale $\mathcal{B}^* = (P_0^*, P_1^*, P_2^*)$.

\subsection{Correction détaillée}

1. \textbf{Démonstration de la liberté de $\mathcal{B}$ :}
   Soient $\lambda_0, \lambda_1, \lambda_2 \in \mathbb{R}$ tels que $\lambda_0 P_0 + \lambda_1 P_1 + \lambda_2 P_2 = 0_E$.
   Cela signifie que pour tout réel $x$, le polynôme s'annule :
   $$\lambda_0 + \lambda_1(x-1) + \lambda_2(x-1)^2 = 0$$
   L'évaluation en $x=1$ donne : $\lambda_0 = 0$.
   En dérivant l'expression polynomiale formellement, nous obtenons :
   $$\lambda_1 + 2\lambda_2(x-1) = 0$$
   L'évaluation en $x=1$ de la dérivée donne : $\lambda_1 = 0$.
   En dérivant une seconde fois : $2\lambda_2 = 0 \implies \lambda_2 = 0$.
   La famille est libre. De cardinal 3 dans un espace de dimension 3, elle constitue une base.

2. \textbf{Recherche de la base duale via les évaluations (formule de Taylor) :}
   Par la formule de Taylor pour les polynômes, tout polynôme $P \in \mathbb{R}_2[X]$ peut s'écrire autour du point $a=1$ :
   $$P(X) = P(1) + P'(1)(X-1) + \frac{P''(1)}{2}(X-1)^2$$
   Substituons avec les vecteurs de notre base :
   $$P(X) = P(1) P_0(X) + P'(1) P_1(X) + \frac{P''(1)}{2} P_2(X)$$
   Par identification avec la décomposition unique dans la base $\mathcal{B}$ : $P = P_0^*(P) P_0 + P_1^*(P) P_1 + P_2^*(P) P_2$, nous identifions les formes coordonnées.

3. \textbf{Vérification de l'action des formes trouvées sur la base :}
   - Soit $P_0^*(P) = P(1)$. On a $P_0^*(P_0) = 1$, $P_0^*(P_1) = 0$, $P_0^*(P_2) = 0$. C'est correct.
   - Soit $P_1^*(P) = P'(1)$. On a $P_1^*(P_0) = 0$, $P_1^*(P_1) = 1$, $P_1^*(P_2) = 2 \times 0 = 0$. C'est correct.
   - Soit $P_2^*(P) = \frac{1}{2}P''(1)$. On a $P_2^*(P_0) = 0$, $P_2^*(P_1) = 0$, $P_2^*(P_2) = \frac{1}{2} \times 2 = 1$. C'est correct.

\textbf{Conclusion :}
La base duale $(P_0^*, P_1^*, P_2^*)$ est formellement caractérisée par les applications : $P \mapsto P(1)$, $P \mapsto P'(1)$ et $P \mapsto \frac{P''(1)}{2}$.

\section{Exercice 4: Intersection d'hyperplans indépendants (Difficulté $\star \star \star$)}

\subsection{Énoncé}
Soit $E$ un espace vectoriel de dimension $n \ge 2$. Soient $H_1$ et $H_2$ deux hyperplans de $E$ dictés par les formes linéaires $\phi_1, \phi_2 \in E^*$. On suppose que $H_1 \neq H_2$. Démontrer avec la plus grande rigueur que la dimension de l'intersection $H_1 \cap H_2$ est égale à $n-2$.

\subsection{Correction détaillée}

1. \textbf{Caractérisation de l'indépendance des formes :}
   Les hyperplans sont donnés par $H_1 = \ker \phi_1$ et $H_2 = \ker \phi_2$.
   Puisque les hyperplans sont distincts ($H_1 \neq H_2$), les formes linéaires $\phi_1$ et $\phi_2$ ne sont pas colinéaires. La famille $(\phi_1, \phi_2)$ est donc une famille libre dans l'espace dual $E^*$.

2. \textbf{Construction de l'application conjointe :}
   Définissons l'application linéaire $\Phi : E \to \mathbb{K}^2$ par $\Phi(x) = (\phi_1(x), \phi_2(x))$.
   La linéarité de $\Phi$ découle directement de la linéarité de ses composantes.

3. \textbf{Détermination du noyau de l'application conjointe :}
   Le noyau de $\Phi$ est l'ensemble des vecteurs $x \in E$ tels que $\Phi(x) = (0, 0)$.
   Cela signifie $\phi_1(x) = 0$ et $\phi_2(x) = 0$.
   Donc $x \in \ker \phi_1 \cap \ker \phi_2$.
   Par conséquent, $\ker \Phi = H_1 \cap H_2$.

4. \textbf{Détermination du rang de l'application conjointe :}
   L'image $\text{Im }\Phi$ est un sous-espace vectoriel de $\mathbb{K}^2$.
   Supposons par l'absurde que $\text{Im }\Phi$ ne soit pas $\mathbb{K}^2$ tout entier. Alors sa dimension serait 0 ou 1.
   Si $\dim(\text{Im }\Phi) = 1$, il existerait $(a, b) \in \mathbb{K}^2 \setminus \{(0,0)\}$ tel que pour tout $x \in E$, l'image soit orthogonale à $(a,b)$ au sens du produit scalaire canonique, soit $a\phi_1(x) + b\phi_2(x) = 0$.
   Cela impliquerait que la combinaison linéaire $a\phi_1 + b\phi_2 = 0_{E^*}$, ce qui contredit formellement la liberté de la famille $(\phi_1, \phi_2)$.
   L'hypothèse est donc absurde. $\Phi$ est surjective, et $\text{rg}(\Phi) = 2$.

5. \textbf{Application du théorème du rang :}
   Appliquons le théorème du rang à l'application $\Phi$ :
   $$\dim(E) = \dim(\ker \Phi) + \text{rg}(\Phi)$$
   En substituant les résultats établis :
   $$n = \dim(H_1 \cap H_2) + 2$$
   $$\dim(H_1 \cap H_2) = n - 2$$

\textbf{Conclusion :}
L'intersection de deux hyperplans distincts dans un espace de dimension $n$ crée structurellement un sous-espace de dimension $n-2$.

\section{Exercice 5: Orthogonal d'un hyperplan (Difficulté $\star \star \star$)}

\subsection{Énoncé}
Soit $E$ un espace vectoriel sur $\mathbb{K}$ de dimension $n$. Soit $H$ un hyperplan de $E$, et $\phi \in E^*$ une forme linéaire non nulle telle que $H = \ker \phi$. Démontrer que l'orthogonal de $H$ dans le dual, noté $H^\perp$, est la droite vectorielle engendrée par $\phi$, c'est-à-dire $H^\perp = \text{Vect}(\phi)$.

\subsection{Correction détaillée}

1. \textbf{Preuve de l'inclusion $\text{Vect}(\phi) \subseteq H^\perp$ :}
   Soit $\psi \in \text{Vect}(\phi)$. Il existe un scalaire $\lambda \in \mathbb{K}$ tel que $\psi = \lambda \phi$.
   Pour démontrer que $\psi \in H^\perp$, il faut vérifier que $\psi$ annule tout vecteur de $H$.
   Soit $x \in H$. Puisque $H = \ker \phi$, nous savons que $\phi(x) = 0$.
   Évaluons $\psi(x)$ :
   $$\psi(x) = (\lambda \phi)(x) = \lambda \phi(x) = \lambda \times 0 = 0$$
   L'application $\psi$ s'annule sur tout $H$, donc formellement $\psi \in H^\perp$.
   Cette première inclusion est validée.

2. \textbf{Évaluation de la dimension par l'isomorphisme canonique :}
   Nous savons, d'après les propriétés fondamentales des orthogonaux en dimension finie, que :
   $$\dim(H) + \dim(H^\perp) = \dim(E)$$
   Puisque $H$ est un hyperplan par hypothèse, le théorème de dimension stipule que $\dim(H) = n - 1$.
   En substituant :
   $$(n - 1) + \dim(H^\perp) = n$$
   D'où $\dim(H^\perp) = 1$.

3. \textbf{Conclusion par argument de dimension :}
   Nous avons établi que $\text{Vect}(\phi)$ est un sous-espace inclus dans $H^\perp$.
   De plus, comme $\phi \neq 0_{E^*}$, la droite $\text{Vect}(\phi)$ possède une dimension égale à 1.
   Puisque $\dim(H^\perp) = 1$, l'inclusion d'un sous-espace de dimension 1 dans un espace de dimension 1 implique l'égalité stricte des ensembles.

\textbf{Conclusion :}
Nous avons $H^\perp = \text{Vect}(\phi)$. L'espace des formes annihilant un hyperplan est structurellement unidimensionnel, engendré par l'équation même de l'hyperplan.

\section{Exercice 6: Théorème d'interpolation de Lagrange vu par la dualité (Difficulté $\star \star \star \star$)}

\subsection{Énoncé}
Soit $E = \mathbb{R}_{n-1}[X]$. Soient $x_1, \dots, x_n$ des nombres réels deux à deux distincts. On définit les formes linéaires d'évaluation $\phi_i : P \mapsto P(x_i)$ pour $i \in \{1, \dots, n\}$. Démontrer que la famille $(\phi_1, \dots, \phi_n)$ est une base de l'espace dual $E^*$.

\subsection{Correction détaillée}

1. \textbf{Dimension de l'espace de travail :}
   L'espace $E = \mathbb{R}_{n-1}[X]$ est l'espace des polynômes de degré strictement inférieur à $n$. Sa dimension canonique est $n$.
   Son dual $E^*$ est par conséquent de dimension $n$.
   La famille de formes $(\phi_1, \dots, \phi_n)$ possède $n$ éléments. Pour prouver qu'il s'agit d'une base, il est suffisant et nécessaire de démontrer que la famille est libre.

2. \textbf{Démonstration de la liberté de la famille de formes :}
   Soient $\lambda_1, \dots, \lambda_n \in \mathbb{R}$ tels que la combinaison linéaire des formes s'annule :
   $$\sum_{i=1}^n \lambda_i \phi_i = 0_{E^*}$$
   Cette relation stipule que l'application combinée renvoie zéro pour *tout* polynôme de $E$. Soit pour tout $P \in \mathbb{R}_{n-1}[X]$ :
   $$\sum_{i=1}^n \lambda_i P(x_i) = 0$$

3. \textbf{Utilisation des polynômes de Lagrange comme "fonctions test" :}
   Pour démontrer que chaque scalaire $\lambda_k$ est nul, nous devons trouver un polynôme qui isole l'évaluation $x_k$ et annule toutes les autres. C'est l'essence des polynômes de Lagrange.
   Pour un indice $k$ fixé, construisons le polynôme $L_k \in \mathbb{R}_{n-1}[X]$ :
   $$L_k(X) = \prod_{\substack{j=1 \\ j \neq k}}^n \frac{X - x_j}{x_k - x_j}$$
   Le degré de $L_k$ est exactement $n-1$, il appartient donc bien à l'espace primal $E$.
   Ses propriétés d'évaluation sont par construction :
   - $L_k(x_k) = 1$
   - $L_k(x_j) = 0$ pour tout $j \neq k$.

4. \textbf{Évaluation de l'identité duale :}
   Appliquons la relation d'annulation sur le polynôme $L_k$ :
   $$\sum_{i=1}^n \lambda_i L_k(x_i) = 0$$
   Par la propriété d'annulation des points distincts, seul le terme d'indice $k$ survit dans la somme :
   $$\lambda_k L_k(x_k) = 0$$
   $$\lambda_k \times 1 = 0 \implies \lambda_k = 0$$
   Ce raisonnement étant symétriquement valide pour tout $k \in \{1, \dots, n\}$, nous déduisons que tous les coefficients $\lambda_i$ sont nuls.

\textbf{Conclusion :}
La famille $(\phi_1, \dots, \phi_n)$ est formellement libre de cardinal égal à la dimension de l'espace. Elle constitue la base duale associée à la base primale des polynômes de Lagrange $(L_1, \dots, L_n)$.

\section{Exercice 7: Trace d'une matrice et forme linéaire (Difficulté $\star \star \star \star$)}

\subsection{Énoncé}
Soit l'espace vectoriel $E = \mathcal{M}_n(\mathbb{R})$. L'application trace $\text{Tr} : M \mapsto \sum_{i=1}^n m_{i,i}$ est une forme linéaire. Démontrer que tout hyperplan de $\mathcal{M}_n(\mathbb{R})$ contient au moins une matrice inversible (pour $n \ge 2$).

\subsection{Correction détaillée}

1. \textbf{Caractérisation formelle de l'hyperplan :}
   Soit $H$ un hyperplan de $\mathcal{M}_n(\mathbb{R})$. Par définition, $H$ est le noyau d'une forme linéaire non nulle $\phi : \mathcal{M}_n(\mathbb{R}) \to \mathbb{R}$.

2. \textbf{Lemme de représentation matricielle (Produit scalaire de Frobenius) :}
   Toute forme linéaire sur $\mathcal{M}_n(\mathbb{R})$ peut s'exprimer par l'intermédiaire de la trace. Il existe une unique matrice $A \in \mathcal{M}_n(\mathbb{R})$, non nulle, telle que pour toute matrice $M$, $\phi(M) = \text{Tr}(A^T M)$.
   Ainsi, $M \in H \iff \text{Tr}(A^T M) = 0$.

3. \textbf{Preuve de l'existence par l'absurde ou par construction :}
   Nous voulons montrer qu'il existe $M$ inversible telle que $\text{Tr}(A^T M) = 0$.
   Considérons le polynôme caractéristique d'une variable matrice. Plus directement, prenons un ensemble fini de matrices inversibles. La condition d'être non inversible (déterminant nul) définit une hypersurface algébrique dans l'espace des matrices (de dimension $n^2-1$), tandis que l'hyperplan vectoriel est de dimension $n^2-1$.
   Puisque l'hyperplan $H$ est un espace vectoriel sur $\mathbb{R}$, il est non borné et connexe. Le sous-ensemble des matrices singulières $S = \{M \in \mathcal{M}_n(\mathbb{R}) \mid \det(M) = 0\}$ est une hypersurface de degré $n$.
   L'intersection $H \cap S$ est un sous-ensemble algébrique propre de $H$. Un polynôme non nul ne peut s'annuler sur un espace vectoriel entier. La seule exception serait si le polynôme déterminant s'annulait identiquement sur l'hyperplan, ce qui est impossible (un espace vectoriel de dimension $n^2-1 > n^2-n$ pour $n \ge 2$ ne peut être totalement contenu dans l'hypersurface de déterminant nul).
   Il existe donc de manière dense des matrices dans $H$ qui n'appartiennent pas à $S$.

\textbf{Conclusion :}
Tout hyperplan vectoriel de l'espace des matrices (dimension $n^2$) est trop vaste pour être confiné au cône des matrices singulières. Il coupe nécessairement le groupe linéaire $GL_n(\mathbb{R})$.

\section{Exercice 8: Bidulalité et annulateur d'annulateur (Difficulté $\star \star \star \star \star$)}

\subsection{Énoncé}
Soit $E$ un espace vectoriel de dimension finie. Soit $F$ un sous-espace vectoriel de $E$. En utilisant l'isomorphisme canonique $\Psi : E \to E^{**}$, démontrer formellement la propriété d'involution orthogonale : $(F^\perp)^\circ = F$, où l'orthogonalité est prise dans le sens des espaces duaux.

\subsection{Correction détaillée}

1. \textbf{Définition rigoureuse des espaces d'annulateurs :}
   Rappelons les définitions :
   - $F^\perp = \{ \phi \in E^* \mid \forall x \in F, \phi(x) = 0 \}$ (Annulateur dans $E^*$ du sous-espace $F \subset E$).
   - Pour un sous-espace $G \subset E^*$, $G^\circ = \{ x \in E \mid \forall \phi \in G, \Psi(x)(\phi) = 0 \} = \{ x \in E \mid \forall \phi \in G, \phi(x) = 0 \}$.
   Le but est de montrer $(F^\perp)^\circ = F$.

2. \textbf{Démonstration de l'inclusion directe $F \subseteq (F^\perp)^\circ$ :}
   Soit $x \in F$.
   Pour prouver que $x \in (F^\perp)^\circ$, il faut vérifier que pour toute forme $\phi \in F^\perp$, l'évaluation donne $\phi(x) = 0$.
   Or, si $\phi \in F^\perp$, par définition même, la forme $\phi$ annule tous les vecteurs de $F$. Comme $x \in F$, on a bien $\phi(x) = 0$.
   Ceci prouve que l'action est réciproque. Donc $x \in (F^\perp)^\circ$.
   L'inclusion $F \subseteq (F^\perp)^\circ$ est établie.

3. \textbf{Démonstration de l'égalité par le théorème des dimensions :}
   Appliquons les théorèmes de dimensions des orthogonaux en dimension finie.
   Nous savons que $\dim(F^\perp) = \dim(E) - \dim(F)$.
   De manière strictement analogue pour le sous-espace $G = F^\perp \subseteq E^*$, on a pour l'annulateur dans le primal :
   $\dim(G^\circ) = \dim(E^*) - \dim(G)$.
   En remplaçant :
   $\dim((F^\perp)^\circ) = \dim(E) - \dim(F^\perp)$
   $\dim((F^\perp)^\circ) = \dim(E) - (\dim(E) - \dim(F))$
   $\dim((F^\perp)^\circ) = \dim(F)$.

4. \textbf{Conclusion :}
   Nous avons une inclusion vectorielle $F \subseteq (F^\perp)^\circ$ entre deux sous-espaces de même dimension finie.
   Cette condition suffit à garantir l'égalité stricte des ensembles.
   $(F^\perp)^\circ = F$.

\section{Exercice 9: Base anté-duale (Difficulté $\star \star \star \star$)}

\subsection{Énoncé}
Soit $E = \mathbb{R}^3$. On donne la famille de formes linéaires $\mathcal{C}^* = (\phi_1, \phi_2, \phi_3)$ définies par :
$\phi_1(x, y, z) = x + y$, $\phi_2(x, y, z) = y + z$, $\phi_3(x, y, z) = x + z$.
Montrer que c'est une base de $E^*$ et trouver la base primale associée $\mathcal{C} = (e_1, e_2, e_3)$ telle que $\mathcal{C}^*$ en soit la base duale.

\subsection{Correction détaillée}

1. \textbf{Vérification du caractère libre :}
   La matrice représentative des formes coordonnées par rapport à la base canonique duale est :
   $M = \begin{pmatrix} 1 & 1 & 0 \\ 0 & 1 & 1 \\ 1 & 0 & 1 \end{pmatrix}$
   Calculons le déterminant par développement selon la première ligne :
   $\det(M) = 1(1 \times 1 - 1 \times 0) - 1(0 \times 1 - 1 \times 1) + 0 = 1 - (-1) = 2 \neq 0$.
   La famille est libre de cardinal 3, c'est une base de $E^*$.

2. \textbf{Recherche de la base primale :}
   Nous cherchons les vecteurs $e_i = (x_i, y_i, z_i)$ tels que $\phi_j(e_i) = \delta_{i,j}$.
   Pour $e_1 = (x_1, y_1, z_1)$ :
   - $\phi_1(e_1) = x_1 + y_1 = 1$
   - $\phi_2(e_1) = y_1 + z_1 = 0 \implies z_1 = -y_1$
   - $\phi_3(e_1) = x_1 + z_1 = 0 \implies x_1 = -z_1 = y_1$
   En remplaçant $x_1$ et $y_1$ dans la première : $y_1 + y_1 = 1 \implies y_1 = \frac{1}{2}$.
   Donc $x_1 = \frac{1}{2}$ et $z_1 = -\frac{1}{2}$. Le vecteur est $e_1 = (\frac{1}{2}, \frac{1}{2}, -\frac{1}{2})$.

3. \textbf{Calcul de $e_2$ et $e_3$ de façon identique :}
   - Pour $e_2$ : $\phi_1(e_2) = 0 \implies y = -x$, $\phi_2(e_2) = 1 \implies y + z = 1$, $\phi_3(e_2) = 0 \implies x = -z$.
     D'où $y = z = \frac{1}{2}$ et $x = -\frac{1}{2}$. $e_2 = (-\frac{1}{2}, \frac{1}{2}, \frac{1}{2})$.
   - Pour $e_3$ : $\phi_1(e_3) = 0 \implies y = -x$, $\phi_2(e_3) = 0 \implies z = -y$, $\phi_3(e_3) = 1 \implies x + z = 1$.
     D'où $x = z = \frac{1}{2}$ et $y = -\frac{1}{2}$. $e_3 = (\frac{1}{2}, -\frac{1}{2}, \frac{1}{2})$.

\textbf{Conclusion :}
La base anté-duale est explicite.

\section{Exercice 10: Séparabilité de classes par un hyperplan (Difficulté $\star \star \star \star \star$)}

\subsection{Énoncé}
Dans $E = \mathbb{R}^n$, soient $A$ et $B$ deux ensembles convexes compacts disjoints non vides. Un théorème de séparation stricte garantit l'existence d'une forme linéaire $\phi \in E^*$ et d'un scalaire $\alpha \in \mathbb{R}$ tels que pour tout $a \in A, \phi(a) < \alpha$ et pour tout $b \in B, \phi(b) > \alpha$.
Montrer que l'hyperplan vectoriel $H = \ker \phi$ permet de définir géométriquement deux demi-espaces stricts séparant les ensembles.

\subsection{Correction détaillée}

1. \textbf{Le rôle fondamental de l'hyperplan affine :}
   L'équation de séparation $\phi(x) = \alpha$ définit un sous-espace affine.
   Puisque $\phi$ est une forme linéaire sur $\mathbb{R}^n$, elle s'écrit $\phi(x) = \langle w, x \rangle$ par le théorème de Riesz (produit scalaire). Le vecteur $w$ est le vecteur normal à l'hyperplan.
   L'hyperplan affine séparateur est $H_\alpha = \{ x \in E \mid \phi(x) = \alpha \}$.
   L'hyperplan vectoriel directeur est $H = \ker \phi = \{ x \in E \mid \phi(x) = 0 \}$.

2. \textbf{Création des zones de classification (le perceptron) :}
   L'espace $E$ est partitionné par $H_\alpha$ en trois parties disjointes :
   - $H_\alpha$ lui-même.
   - $E^- = \{ x \in E \mid \phi(x) < \alpha \}$ (le demi-espace ouvert inférieur).
   - $E^+ = \{ x \in E \mid \phi(x) > \alpha \}$ (le demi-espace ouvert supérieur).
   L'hypothèse du théorème stipule exactement que l'ensemble compact $A$ est entièrement inclus dans $E^-$, et l'ensemble $B$ est entièrement inclus dans $E^+$.

3. \textbf{Caractère topologique strict :}
   Parce que les inégalités sont strictes, la frontière de décision $H_\alpha$ (le séparateur de marge) ne touche ni l'ensemble $A$ ni l'ensemble $B$.
   La distance (la marge) de $H_\alpha$ à $A$ et $B$ est strictement positive en raison de la compacité des ensembles (la forme continue $\phi$ atteint ses bornes sur un compact).

\textbf{Conclusion :}
C'est la formalisation mathématique absolue de la capacité d'un neurone artificiel (perceptron) à discriminer linéairement deux clusters de données distincts dans un espace latent.


\chapter{Travaux Pratiques en Python}
\section{TP 1: Implémentation algorithmique du Dual}

\subsection{Objectif mathématique}
Ce TP exige la modélisation en Python pur de la dualité, validant l'action d'évaluation sans faire appel à numpy.

\subsection{Code Python}
\begin{lstlisting}

def forme_lineaire_eval(coefficients_duaux, vecteur_primal):
    """
    Évalue la forme linéaire phi(x) définie par la matrice des coefficients duaux.
    phi(x) = a1*x1 + a2*x2 + ... + an*xn
    """
    if len(coefficients_duaux) != len(vecteur_primal):
        raise ValueError("L'instrument de mesure et l'objet doivent appartenir à la même dimension.")

    somme_evaluation = 0.0
    for i in range(len(vecteur_primal)):
        somme_evaluation += coefficients_duaux[i] * vecteur_primal[i]

    return somme_evaluation

# Verification rigoureuse
coef_phi = [2.0, -1.5, 3.0]
vecteur = [1.0, 2.0, -1.0]

resultat = forme_lineaire_eval(coef_phi, vecteur)
# 2(1) - 1.5(2) + 3(-1) = 2 - 3 - 3 = -4
assert resultat == -4.0
print("Évaluation duale validée mathématiquement.")
\end{lstlisting}
\section{TP 2: Calcul des composantes d'une base duale}

\subsection{Objectif mathématique}
Calcul matriciel d'inversion pour trouver les formes cordonnées par la méthode de pivot algébrique pure.

\subsection{Code Python}
\begin{lstlisting}

def produit_scalaire(u, v):
    return sum(ui * vi for ui, vi in zip(u, v))

def soustraire_vecteurs(u, v, facteur=1.0):
    return [ui - facteur * vi for ui, vi in zip(u, v)]

def base_orthonormale_gram_schmidt(base):
    # Implementation fondamentale pour extraire la dualite dans un cadre euclidien
    # Pour un repere ON, primal et dual se superposent via le vecteur normal
    base_onb = []
    for vecteur in base:
        u = list(vecteur)
        for base_u in base_onb:
            proj_coef = produit_scalaire(vecteur, base_u) / produit_scalaire(base_u, base_u)
            u = soustraire_vecteurs(u, base_u, proj_coef)

        # Normalisation mathematique stricte
        norme_u = produit_scalaire(u, u) ** 0.5
        if norme_u > 1e-9:
            u_norme = [ui / norme_u for ui in u]
            base_onb.append(u_norme)
        else:
            raise ValueError("La famille originelle n'était pas libre.")
    return base_onb

# Validation
base_canonique = [[1,0,0], [0,1,0], [0,0,1]]
ortho = base_orthonormale_gram_schmidt(base_canonique)
assert produit_scalaire(ortho[0], ortho[1]) == 0.0
print("Extraction algorithmique des composantes primales terminée.")
\end{lstlisting}
\section{TP 3: Séparateur de données (Hyperplan)}

\subsection{Objectif mathématique}
Modéliser un hyperplan vectoriel et diviser l'espace en deux demi-espaces fermés via le noyau.

\subsection{Code Python}
\begin{lstlisting}

def classifier_par_hyperplan(vecteur_normal, point):
    """
    Évalue si un point est au-dessus (1), sur (0) ou en-dessous (-1) de l'hyperplan.
    """
    evaluation = 0.0
    for i in range(len(vecteur_normal)):
        evaluation += vecteur_normal[i] * point[i]

    if evaluation > 0:
        return 1
    elif evaluation < 0:
        return -1
    else:
        return 0  # Appartient au noyau, donc à l'hyperplan

# Test canonique de la frontiere
w_forme_lineaire = [1, -1] # Hyperplan: x - y = 0
assert classifier_par_hyperplan(w_forme_lineaire, [2, 1]) == 1   # 2-1 = 1 > 0
assert classifier_par_hyperplan(w_forme_lineaire, [1, 2]) == -1  # 1-2 = -1 < 0
assert classifier_par_hyperplan(w_forme_lineaire, [3, 3]) == 0   # 3-3 = 0
print("Topologie des demi-espaces confirmée algorithmiquement.")
\end{lstlisting}
\section{TP 4: Validation de l'intersection de sous-espaces duaux}

\subsection{Objectif mathématique}
Détecter que l'intersection des noyaux de deux formes linéaires non colinéaires induit une chute de dimension de 2 dans le primal.

\subsection{Code Python}
\begin{lstlisting}

def dependance_lineaire(phi1, phi2):
    """
    Vérifie si deux formes linéaires (listes de coefficients) sont colinéaires.
    """
    n = len(phi1)
    if n == 0: return True

    # Trouver la premiere composante non nulle pour fixer le rapport
    ratio = None
    for i in range(n):
        if phi1[i] != 0:
            if phi2[i] == 0:
                return False
            ratio = phi1[i] / phi2[i]
            break

    if ratio is None:
        return phi2 == [0]*n # Les deux sont nulles

    # Verification universelle du ratio
    for i in range(n):
        if abs(phi1[i] - ratio * phi2[i]) > 1e-9:
            return False

    return True

# Validation
assert dependance_lineaire([1, 2, 3], [2, 4, 6]) == True
assert dependance_lineaire([1, 2, 3], [1, 4, 6]) == False
print("Indépendance algébrique des formes vérifiée avec succès.")
\end{lstlisting}
\section{TP 5: Approximation du Bidual}

\subsection{Objectif mathématique}
Simuler algorithmiquement l'application d'isomorphisme canonique vers le Bidual: x -> \text{ev}_x.

\subsection{Code Python}
\begin{lstlisting}

def \text{ev}_x(vecteur_x):
    """
    Construit la fonction d'évaluation canonique (le bidual) pour un vecteur donné.
    Retourne une fonction qui accepte une forme linéaire.
    """
    def isomorphisme(forme_phi):
        # Application stricte de phi(x)
        resultat = 0.0
        for i in range(len(vecteur_x)):
            resultat += forme_phi[i] * vecteur_x[i]
        return resultat
    return isomorphisme

# Validation structurelle
vecteur_primal = [5, -2]
bidual_instance = \text{ev}_x(vecteur_primal)

# La forme lineaire phi(x,y) = 3x + y
forme_test = [3, 1]

# Evaluation via le bidual: Psi(x)(phi)
evaluation_biduale = bidual_instance(forme_test)

# Evaluation directe: phi(x) = 3(5) + 1(-2) = 13
assert evaluation_biduale == 13.0
print("Isomorphisme du Bidual codé et validé fonctionnellement.")
\end{lstlisting}

\end{document}